\scanpage{45}

prenant le det du faisceau descendu, on trouve à partir d'un faisceau \unsure{sur ample|très ample|rel. ample} sur $X/Y$, un faisceau ample sur $f_*(X/T)/f_*(Y/T)$ \deleted{etc}]. Kif-kif pour propre.

\uline{Proposition} Soit $g:X_{/T}\to Y_{/T}$ un morphisme, et soit $\underline L$ un faisceau inversible sur $X$ \uline{ample} \deleted{pour} (resp. très ample) pour $g$. Alors $\det f_*(\underline L)$ est un faisceau ample (resp. très ample) pour $f_*(X/T)\longrightarrow f_*(Y/T)$

Démonstration. \textcircled{2} on peut supposer $Y$ affine, \inserted{(Sous la réserve faite)} l'amplitude \deleted{de \unsure{très-amplitude|très amplitude}} étant une notion locale. $^{(1)}$ On \deleted{\unsure{voit par un|voit par ce} \unsure{que|qui} en \unsure{endroit}} \inserted{est ramené au cas ``ample''.} De plus, on est ramené à prouver ceci : Si $X=P^m_Y$, et $\underline L$ le faisceau fondamental dessus, \deleted{alors} $f_*(\underline L)/T)$ \deleted{est} définit une immersion projective fermée de $f_*(P^m_Y)$ par rapport à $f_*(Y/T)$. Mais utilisant la commutation de $f_*$ avec les produits, on est ramené au cas où $Y=T$. De plus, la question étant locale, on peut supposer $S$ affine $=\operatorname{Spec}(A)$, $T=\operatorname{Spec}(B)$, $B$ étant libre de type fini (rang $d$) sur $A$.

\margin{$f_*(\unsure{X/T|U/T|\mathrm{Ker}})\longrightarrow f_*(Y/T)$ \qquad $S\qquad T$}

% slanted
\margin{\deleted{N.B. \unsure{théorème|théor.|Résultat} s'améliore si on introduit sur $Y$ le faisceau $g_*(\underline L)$ (par hypo fini ou quasicoh.) \uline{pourvu} \uline{que} (i) les \deleted{\unsure{foncteurs}} ouverts \unsure{d'imm.|d'une} projectives marchent bien sans réserve aucune. (ii) On le prouve pour un \uline{tel} faisceau quasi-cohérent}}

On est ramené \deleted{au} à deux vérifications :

A) Détermination de la grassmannienne. et d'un faisceau ample dessus, \deleted{\unsure{id.|i.e.}} d'une immersion projective de la dite.

% archive sheet number: 44
